\section{Exact events and their relation spaces}
\label{p13:exact:section}

The optimal remainder is smaller than \(N^{-1}\), the cost of deleting
all irregular windows. We therefore keep every event with nonzero
probability from the outset. Its effective square-class space records
all local relations once, and its actual probability is the correct
centering parameter.

Let \(\mathcal I_N=\{x\in I_N:p_x>0\}\), where \(p_x=\E J_x\).
An event outside \(\mathcal I_N\) is identically zero and can be
removed without changing \(Z_N\). For every \(x\in I_N\), let
\(S_x\) be the image of the even selections of
\(V_x=\{x-1,\ldots,x+L-1\}\) in
\(\mathbb Q^\times/\mathbb Q^{\times2}\), regarded as a vector
space over \(\F_2\). Write \(d_x=\dim S_x\).
The even-selection space has dimension $L$. Its kernel consists of
local even square relations. The following bound makes the possible
effective ranks explicit.

Lemma~\ref{p1:limits:local-rigidity} in the companion proves that
the local even square-relation code has dimension at most one in the
critical band. Its proof combines a short-interval square-product bound
with the binary Griesmer bound.

For an admissible window the affine prescription is consistent. Its
probability is $2^{-d_x}$, as the character expansion in
Section~\ref{p1:alg:section} also shows. The lemma therefore gives
\begin{equation}
 d_x\in\{L-1,L\},\qquad p_x=2^{-d_x}\in\{p,2p\},
 \qquad p_*:=\max_{x\in\mathcal I_N}p_x\ll_C N^{-1}.
 \label{p13:exact:marginals}
\end{equation}
The assertion about \(d_x\) holds for all sufficiently large \(N\)
in a fixed critical band. A local square relation may instead
contradict the run prescription; that is precisely the case
already removed by \(p_x=0\).

For a finite \(I\subseteq\mathcal I_N\), put
\begin{align}
 \rho_{\mathrm{rel}}(I)
  &=\sum_{x\in I}d_x-\dim\sum_{x\in I}S_x,
       & d_{xy}&=\dim(S_x\cap S_y),\label{p13:exact:rank}\\
 q(I)&=\E\prod_{x\in I}J_x,
       & w(I)&=2^{\rho_{\mathrm{rel}}(I)}
                         \prod_{x\in I}p_x,\label{p13:exact:weights}\\
 c(I)&=\E\prod_{x\in I}(J_x-p_x).
 \label{p13:exact:center}
\end{align}
The empty products are one. We have \(q(I)=w(I)\) when the
affine prescriptions are jointly compatible, and \(q(I)=0\)
otherwise. Write $G_N=\{x\in I_N:d_x=L\}$ for the regular set.
On $G_N$, these definitions agree
with the homogeneous-marginal notation used in the companion.
Outside \(G_N\), \(\rho_{\mathrm{rel}}\) counts dependencies
between effective spaces; it does not count local relations a
second time.

The notions of directness, connectedness, and primitive family
refer to the direct-sum decomposition of the effective spaces
proved below. A family is direct when its
effective spaces are in direct sum, and connected when it has just one
canonical direct component. In particular, a primitive
family has a nonzero relation and every proper subfamily is
direct. Every nonzero relation of a primitive family has a
nonzero projection in each of its blocks. A connected family
need not have a relation with this property on its entire
support. This distinction matters in the conditional count below.

For fixed \(r\ge2\), define
\begin{equation}
 A_r=\sum_{\substack{I\subseteq\mathcal I_N\\|I|=r}}|c(I)|,
 \qquad
 K_3^+=\sum_{\substack{|T|=3,\ T\text{ connected}\\q(T)>0}}c(T).
 \label{p13:exact:activities}
\end{equation}
The nonnegativity in the definition of \(K_3^+\) will be proved
in Proposition~\ref{p13:activity:kappa3}.
The exact second and third factorial cumulants satisfy
\begin{align}
 \alpha_N&=2\sum_{|e|=2}c(e)-\sum_xp_x^2,
 \label{p13:exact:alpha}\\
 \kappa_{3,N}
   &=6\sum_{|T|=3}c(T)
       -6\sum_{e=\{x,y\}}(p_x+p_y)c(e)
       +2\sum_xp_x^3.
 \label{p13:exact:kappa3}
\end{align}
These identities follow by expanding the first three factorial
moments. They display the repeated-index terms that an
independent Bernoulli reference will retain exactly.

\paragraph{Fixed sources on effective spaces.}
The fixed-source estimates stated in
Propositions~\ref{p1:fixed:targets} and~\ref{p1:fixed:propagation}
apply on
\(\mathcal I_N\), with the new conditional rank
\[
 \delta_z=\dim\bigl(S_z\cap\textstyle\sum_{x\in I}S_x\bigr)
\]
and weight \(p_z2^{\delta_z}\). To see this, choose in each
even-selection space a linear complement to its local kernel.
The valuation map is injective on that complement and maps it
onto \(S_x\). Every label regularization, prime row, and disjoint
square equation in the proofs is applied to this linear
realization. Restriction to a complement restricts the allowed
binary selections within the same physical coordinates; it does
not enlarge any prime support.
The factor \(p_z\) is at most \(2p\), and the number of sources
and targets is fixed in those propositions. Thus their constants
change by a fixed factor. The same reasoning applies to
homogeneous prescriptions. In particular, for fixed \(s,m\) and
\(\varepsilon>0\),
\begin{equation}
 \sum_{\substack{T\subseteq\mathcal I_N\setminus I\\|T|=m}}
       q(I\cup T)\ll_{s,m,C,\varepsilon}N^\varepsilon q(I)
 \quad (|I|=s),
 \label{p13:exact:propagation}
\end{equation}
and the analogous homogeneous bound holds with \(q\) replaced
by \(w\). If \(q(I)=0\), every summand on the left is zero.
No estimate here is asserted for \(s\) or \(m\) growing with \(N\).

\subsection{Affine structure and centered coefficients}
\label{p1:alg:section}

The arithmetic estimates naturally count relations, whereas the
approximation problem concerns a signed difference of laws.
We introduce the exact identities linking these two objects.
They will be used with a stopping rule, so that large connected
families outside the required nullity range never need to be
summed in absolute value.

Throughout this section $I\subseteq\mathcal I_N$, and
$\rho_{\mathrm{rel}}(I)$ is the effective relative nullity.
Write $b_x=\delta_0+p_x\Delta$. The identities below hold for
heterogeneous marginals. On the regular set $G_N$, one has $p_x=p$ and $d_x=L$.
The same identities therefore cover both regular and irregular windows.

\subsubsection{An exact expansion around independent Bernoulli variables}

Set
\[
 Y_x=J_x-p_x,\qquad c(I)=\E\prod_{x\in I}Y_x,
 \qquad c(\varnothing)=1.
\]

\begin{proposition}[Centered Bernoulli expansion]
\label{p1:alg:centred}
For every \(N\),
\begin{equation}
 \Law(Z_N)=\sum_{I\subseteq\mathcal I_N}
       c(I)\Delta^{*|I|}*\mathop{*}_{x\notin I}b_x .
 \label{p1:alg:centred-identity}
\end{equation}
The complementary convolution products are deterministic probabilities.
The terms of sizes zero and two are
\[
 \mathop{*}_{x}b_x+
 \sum_{|e|=2}c(e)\Delta^{*2}*\mathop{*}_{x\notin e}b_x.
\]
Terms of size one vanish.
\end{proposition}

\begin{proof}
For an indicator,
\(\delta_{J_x}=b_x+Y_x\Delta\). Expand the finite convolution
product over \(\mathcal I_N\) and take expectations. No independence
assumption is used in this identity.
\end{proof}

Each \(s\in S_x\) is a square-class vector. Let
\(\chi_s\) be its character on the prime signs, and let
\(\eta_x(s)\in\{\pm1\}\) be the affine sign prescribed by
the run-start event. The local prescription is consistent on
\(\mathcal I_N\), so its affine character \(\eta_x\) is well defined
on the effective space. Thus
\[
 J_x=p_x\sum_{s\in S_x}\eta_x(s)\chi_s,\qquad \eta_x(0)=1.
\]

\begin{lemma}[Full-support Fourier formula]
\label{p1:alg:fourier}
For nonempty \(I\subseteq\mathcal I_N\),
\begin{equation}
 c(I)=\Bigl(\prod_{x\in I}p_x\Bigr)
   \sum_{\substack{s_x\in S_x\setminus\{0\}\ (x\in I)\\
                         \sum_xs_x=0}}
                 \prod_{x\in I}\eta_x(s_x).
 \label{p1:alg:fourier-identity}
\end{equation}
Consequently,
\[
 |c(I)|\le\Bigl(\prod_{x\in I}p_x\Bigr)(2^{\rho_{\mathrm{rel}}(I)}-1)
           \le\Bigl(\prod_{x\in I}p_x\Bigr)2^{\rho_{\mathrm{rel}}(I)}.
\]
If \(c(I)\ne0\), there is a relative relation meeting
every block of \(I\).
\end{lemma}

\begin{proof}
Subtracting \(p_x\) removes the zero character from each
factor. The expectation of a character product is zero
unless the sum of its indices vanishes. The kernel of
\(\bigoplus_{x\in I}S_x\to\sum_{x\in I}S_x\) has
\(2^{\rho_{\mathrm{rel}}(I)}\) elements. The displayed sum omits its
zero element and possibly others. Every affine
coefficient has modulus one.
\end{proof}

We extend the same formulas for $q(I)$, $c(I)$ and relative nullity
to arbitrary original window families when stating the next lemma.

\begin{lemma}[A common homogeneous envelope]
\label{p2:homogeneous:envelope}
For a nonempty family $I\subseteq I_N$ of original windows, set
\[
 q_0(I)=2^{-\dim(\sum_{x\in I}S_x)},\qquad
 w(I)=\Bigl(\prod_{x\in I}p_x\Bigr)2^{\rho_{\mathrm{rel}}(I)}.
\]
Then
\begin{equation}
 q_0(I)=2^{\rho_{\mathrm{rel}}(I)}\prod_{x\in I}2^{-d_x},
 \qquad w(I),\ q(I),\ |c(I)|\le q_0(I).
 \label{p2:homogeneous:domination}
\end{equation}
These assertions allow locally inconsistent start prescriptions and
zero marginals. In the critical band, for all sufficiently large $N$,
\begin{equation}
 q_0(I)\le 2^{|I|}p^{|I|}2^{\rho_{\mathrm{rel}}(I)}.
 \label{p2:homogeneous:nominal}
\end{equation}
\end{lemma}

\begin{proof}
Rank--nullity gives the displayed factorization. A consistent affine
prescription has the same probability as the corresponding homogeneous
system; an inconsistent one has probability zero. Thus
$p_x\le2^{-d_x}$ and $q(I)\le q_0(I)$, proving the assertions for
$w(I)$ and $q(I)$.
If some $p_x=0$, then $J_x=0$ almost surely and $c(I)=0$.
Otherwise the full-support Fourier formula in
Lemma~\ref{p1:alg:fourier} gives $|c(I)|\le w(I)$.
Finally, local rigidity gives $d_x\ge L-1$ uniformly on $I_N$ in
the critical band, so $2^{-d_x}\le2p$ and the last inequality follows.
\end{proof}

The homogeneous quantity is an actual probability: it prescribes zero
for every relative bit. It can be positive even when an original start
is impossible. Its finite domination is valid at every cardinality;
absorbing the factor $2^{|I|}$ into a constant requires a fixed cardinality.

\subsubsection{The nullity that survives affine compatibility}

Define
\begin{equation}
 \rho_*(I)=\max\{\rho_{\mathrm{rel}}(A):A\subseteq I,\ q(A)>0\},
 \qquad q(\varnothing)=1,\quad\rho_{\mathrm{rel}}(\varnothing)=0 .
 \label{p1:alg:compatible-nullity}
\end{equation}
This parameter records the largest rank loss in a
compatible subfamily. It can be smaller than the
nullity of the whole family.

\begin{lemma}[Compatible-nullity majorant]
\label{p1:alg:compatible}
For every \(I\subseteq\mathcal I_N\),
\begin{equation}
 |c(I)|\le 2^{|I|}\Bigl(\prod_{x\in I}p_x\Bigr)2^{\rho_*(I)},\qquad
 \rho_*(I)\le\rho_{\mathrm{rel}}(I).
 \label{p1:alg:compatible-bound}
\end{equation}
Equality holds in the second inequality when \(q(I)>0\).
\end{lemma}

\begin{proof}
The exact inclusion-exclusion formula is
\[
 c(I)=\sum_{A\subseteq I}(-1)^{|I|-|A|}q(A)
                       \prod_{x\in I\setminus A}p_x.
\]
A compatible term has absolute value
\(\bigl(\prod_{x\in I}p_x\bigr)2^{\rho_{\mathrm{rel}}(A)}\), and an incompatible term is
zero. There are at most \(2^{|I|}\) terms.
Nullity increases under inclusion, because adding the block
\(S_x\) increases the total rank by at most \(d_x\). This proves the other assertions.
\end{proof}

\subsubsection{Canonical direct components}

A partition \(\mathcal P\) of \(I\) is called direct if
the spaces \(V_C=\sum_{x\in C}S_x\), \(C\in\mathcal P\),
are in direct sum. Its finest such partition is the
canonical direct decomposition of \(I\).
We call a nonempty family \emph{connected} if this
partition has only one part.

\begin{lemma}[Existence, uniqueness, and factorization]
\label{p1:alg:components}
Every finite family \(I\) has a unique finest direct
partition \(\mathcal P(I)\). For it,
\begin{equation}
 q(I)=\prod_{C\in\mathcal P(I)}q(C),\qquad
 c(I)=\prod_{C\in\mathcal P(I)}c(C),\qquad
 \rho_{\mathrm{rel}}(I)=\sum_{C\in\mathcal P(I)}\rho_{\mathrm{rel}}(C).
 \label{p1:alg:factorization}
\end{equation}
\end{lemma}

\begin{proof}
If \(\mathcal P,\mathcal Q\) are direct partitions,
their nonempty intersections form a direct partition.
Indeed, inside each \(P\in\mathcal P\), the spaces
\(\sum_{x\in P\cap Q}S_x\) lie in the mutually direct
spaces \(V_Q\), and sum to \(V_P\). The direct sum over
\(P\) proves the claim. Intersecting the finitely many
direct partitions gives the unique finest one.

Characters indexed by mutually direct subspaces of
the finite prime-sign vector space are independent
uniform character families. Events, and centered
variables, depending on different \(V_C\) are therefore
independent. This proves both probability identities;
rank additivity proves the last one.
\end{proof}

Connectedness here is defined by all relative
relations. A minimal dependent family of original
blocks is primitive, but primitive block supports
alone do not determine the canonical components:
an additional relation can join a third block to
an already dependent pair. Exact pair compression
will retain this possibility.
